\section{投资问题：谁控制资产，谁承担风险，谁分享利润}

恒逸石化不是一座单体炼厂，也不是纯聚酯制造商，而是由上市公司、控股子公司、合营联营平台和海外项目共同组成的资产网络。投资者需要回答的不是“公司有没有完整产业链”，而是每个利润池以何种比例并表、经过多少少数股东和债务、最终有多少归属于普通股东。核心判断是，控股股东控制稳定、产业运营经验深，但复杂参控股结构、日常关联交易、对外担保和海外项目融资提高了财务穿透难度；这些因素应压低峰值利润的估值倍数。

截至2025年报，浙江恒逸集团直接持股47.04\%，并通过其控股的杭州恒逸投资持股10.61\%；实际控制人为邱建林。公司披露截至2026-03-31，因可转债转股，恒逸集团与恒逸投资合计持股比例被动稀释至54.71\%。此后可转债仍持续转股，因此本报告不把一季度比例写成7月22日实时比例；但控股权稳定这一方向没有改变。

\begin{exhibitbox}[控制权与主要治理节点]
\noindent\begin{tabularx}{\textwidth}{L{3.4cm}L{2.6cm}L{2.6cm}X}
\toprule
\textbf{节点} & \textbf{已披露状态} & \textbf{投资含义} & \textbf{边界／监控} \\
\midrule
实际控制人 & 邱建林；境内自然人 & 长期产业控制与决策连续性 & 控制人并非上市公司董事长兼总经理 \\
控股股东 & 浙江恒逸集团 & 直接控制与产业协同 & 与上市公司之间独立性依赖治理执行 \\
合计持股 & 截至2026Q1为54.71\% & 控制权稳固 & 转债继续转股使比例动态变化 \\
董事长／总裁 & 邱奕博兼任 & 经营决策集中、执行链短 & 与控股股东存在重要股东及任职联系 \\
董事会专门委员会 & 审计、风险控制、薪酬提名、战略投资与ESG & 形式治理架构完整 & 报告仍需用关联交易、担保和资本配置结果验证有效性 \\
审计机构 & 中兴华，连续服务7年 & 历史口径具有连续性 & 2026H1预告仍未经审计 \\
\bottomrule
\end{tabularx}
\end{exhibitbox}
\sourcenote{HY-OFF-2025AR，公司治理与股东章节；HY-OFF-CB，转股截止2026-06-30。持股比例54.71\%为2026-03-31披露值，不代表2026-07-22实时比例。}

稳定控制权有利于跨周期推进文莱一期、二期和国内一体化项目，却也意味着资本配置高度依赖同一控制体系。对投资者而言，控制稳定不是估值溢价的充分条件；只有当项目回报、融资约束和现金分配同时可验证时，稳定性才转化为每股价值。

\section{资产组织：并表、权益法与需求锚不能混在一起}

文莱一期由恒逸文莱经营，上市公司持有70\%权益，属于核心并表利润池，但30\%净利润归少数股东。浙江逸盛为70\%控制的PTA／PIA平台；海南逸盛为50\%间接持股的合营平台；逸盛新材料为49\%持股的联营平台。国内聚酯资产以控股制造实体为主，但“参控股产能”与合并报表产销量并不一致。若把所有产能视为100\%并表，就会同时高估收入、利润和现金流。

\begin{exhibitbox}[资产与归属关系示意]
\centering
\begin{tikzpicture}[
  node distance=7mm and 12mm,
  parent/.style={draw=navy, fill=blue!5, rounded corners=2pt, text width=3.0cm, align=center, minimum height=1.0cm, font=\small},
  asset/.style={draw=rulegrey, fill=lightgrey, rounded corners=2pt, text width=3.2cm, align=center, minimum height=1.15cm, font=\small},
  equity/.style={draw=riskamber, fill=yellow!6, rounded corners=2pt, text width=3.2cm, align=center, minimum height=1.15cm, font=\small},
  line/.style={-{Stealth[length=2mm]}, draw=steelgrey, thick}
]
\node[parent] (listed) {恒逸石化上市公司\\归母利润与每股价值};
\node[asset, below left=of listed] (brunei) {恒逸文莱\\70\%控股\\炼油、PX、苯};
\node[asset, below=of listed] (domestic) {国内控股制造平台\\PTA／PIA、聚酯、CPL／PA6};
\node[equity, below right=of listed] (associates) {合营／联营平台\\海南逸盛50\%\\逸盛新材料49\%};
\draw[line] (brunei) -- node[left,font=\scriptsize]{并表后扣少数} (listed);
\draw[line] (domestic) -- node[right,font=\scriptsize]{并表} (listed);
\draw[line] (associates) -- node[right,font=\scriptsize]{权益法／合营} (listed);
\end{tikzpicture}
\end{exhibitbox}
\sourcenote{HY-OFF-2025AR，主要子公司及在其他主体中的权益；研究截止2026-07-22。图示用于归属逻辑，不代表完整法律实体清单。}

这张归属图解释了为什么本报告拒绝伪精确SOTP：公开披露没有把各法律实体的税后利润、净债务、内部交易和少数股东完整分配到同一时点。主模型先在集团归母层面做周期情景，再把分部证据用于约束，而不是先虚构每个项目的独立价值再相加。

\section{管理层的产业经验是优势，但不替代资本纪律}

董事长兼总裁邱奕博具有集团投资发展和石化销售经历；副董事长方贤水拥有三十年以上化纤生产管理经验并在多家产业平台任职；常务副总裁王松林长期从事石化化纤技术与经营；副总裁陈连财曾任大型炼化企业管理职务，现任恒逸文莱CEO；董事会秘书兼财务总监郑新刚具有投融资和资本市场经验。管理团队对炼化、聚酯和跨境项目的行业认知，是文莱一期运行和国内链条协同的重要能力来源。

但管理经验无法消除三个约束。第一，大型项目的资本回报受周期和融资成本决定；第二，参控股网络增加内部定价和现金归集复杂度；第三，文莱二期在融资条款未完全披露前，执行能力不能等同于股东回报。投资者应观察项目是否在承诺节点内形成可分配现金，而不是只观察产能是否开工。

\begin{keyinsight}[治理的“所以呢”]
产业经验使恒逸能够管理复杂装置和跨境供应链，但估值溢价必须来自可验证的资本回报。\textbf{行动上，二期项目只有在融资交割、资本金、进度和回报门槛同时透明后，才从零基准信用升级为正期权；管理层履历本身不进入目标价。}
\end{keyinsight}

\section{关联交易：产业协同与穿透难度同时存在}

2025年日常关联交易合计约416.10亿元，涉及PTA、PX、PIA、聚酯产品、原油、能源品、运输和工程管理等。大额交易在一体化产业集团中并不意外，公司披露以市场价格或电力部门价格为定价依据，且报告期不存在控股股东非经营性资金占用。但关联交易规模较大，会使合并分部收入、内部消化、外销量和利润分配更难仅凭表面产能判断。

其中，向逸盛新材料采购PTA约131.03亿元，向香港逸天采购原油约124.40亿元，向绍兴恒鸣采购聚酯产品约90.13亿元，均属于高影响交易。关联交易的投资风险不是简单认定价格不公允，而是：当产品价格快速波动时，关联实体之间的定价时点、结算方式、信用期和库存归属可能改变利润与现金流在各主体之间的表现。

\begin{exhibitbox}[高影响关联交易与研究处理]
\small
\noindent\begin{tabularx}{\textwidth}{L{2.7cm}L{2.25cm}R{2.0cm}L{2.2cm}X}
\toprule
\textbf{交易方} & \textbf{主要内容} & \textbf{2025金额} & \textbf{披露定价} & \textbf{研究处理} \\
\midrule
逸盛新材料 & 采购PTA & 约131.03亿元 & 市场价 & 关注内部消化、信用期与权益法归属 \\
香港逸天 & 采购原油 & 约124.40亿元 & 市场价 & 关注原油贴水、库存与跨境结算 \\
绍兴恒鸣 & 采购聚酯产品 & 约90.13亿元 & 市场价 & 不把采购规模写成终端需求增长 \\
海南逸盛 & 采购PTA & 约21.40亿元 & 市场价 & 关注合营平台与上市公司利润分配 \\
全部日常关联交易 & 多品种与服务 & 约416.10亿元 & 市场价／规定电价 & 不直接作负面判断；提高穿透与现金流监控权重 \\
\bottomrule
\end{tabularx}
\end{exhibitbox}
\sourcenote{HY-OFF-2025AR，重大关联交易表，金额由万元换算为亿元，财务截止2025-12-31。}

因此，本报告在价值链中优先使用集团审计分部、公司产销量和子公司利润，而不以关联采购销售额推导外部份额。半年报若能提供分部与关联交易的变化，有助于判断2026H1高利润究竟来自外部价差、内部库存移动还是权益法贡献。

\section{担保与资本约束：没有违规，不等于没有风险}

公司披露2025年末实际担保余额约279.01亿元，占公司净资产114.41\%；其中直接或间接为资产负债率超过70\%的被担保对象提供的债务担保余额约110.67亿元。公司报告期无违规对外担保，也没有证据表明已发生代偿责任；但担保规模意味着信用风险与融资风险在集团内部高度联动。

公司及子公司还为银团向恒逸文莱提供17.5亿美元或等值境外人民币项目贷款担保。该担保支持一期项目融资与运营，却也表明文莱风险不能只在子公司层面隔离。一旦裂解价差、汇率、物流或项目负荷恶化，上市公司信用与再融资空间可能同时受到影响。

\begin{exhibitbox}[担保与治理风险仪表盘]
\noindent\begin{tabularx}{\textwidth}{L{4.1cm}R{2.2cm}L{2.3cm}X}
\toprule
\textbf{指标} & \textbf{2025年末} & \textbf{风险含义} & \textbf{监控动作} \\
\midrule
实际担保余额 & 279.01亿元 & 超过净资产，集团信用联动高 & 半年报更新余额与到期结构 \\
担保余额／净资产 & 114.41\% & 资本缓冲对被担保主体经营敏感 & 与现金、授信和短债一起看 \\
高负债被担保对象余额 & 110.67亿元 & 尾部代偿风险更集中 & 观察被担保主体现金流与展期 \\
股东及实控人担保余额 & 0 & 未向控制人提供担保 & 保持为治理底线 \\
恒逸文莱项目贷款担保 & 17.5亿美元或等值境外人民币 & 海外项目风险穿透上市公司 & 跟踪贷款、汇率、负荷和偿债覆盖 \\
违规担保／非经营性占用 & 无／无 & 没有已披露治理红线 & 不能因此忽略正常担保规模 \\
\bottomrule
\end{tabularx}
\end{exhibitbox}
\sourcenote{HY-OFF-2025AR，重大担保、资金占用与违规担保章节，财务截止2025-12-31。金额由万元／美元口径换算并保留“约”。}

担保风险的估值含义是，它更适合通过较低PE和资本成本体现，而不是从权益价值中机械减去全部担保余额。担保并非当期债务，也未发生代偿；但它限制了公司在周期下行时的融资弹性，是核心公允区间不向牛情景靠拢的重要原因。

\section{治理结论：看资本结果，不只看制度表述}

公司披露董事会治理、独立性和内部控制架构符合规范，2025年没有重大诉讼、处罚、违规担保、非经营性资金占用或非标准审计意见。与此同时，大额关联交易、担保、可转债、员工持股计划和文莱二期共同构成复杂资本结构。治理判断因此不能停在“制度是否存在”，而应继续看资本配置是否提高每股现金价值。

\begin{riskbox}[治理失效条件]
若关联交易应收／预付款显著增加、担保余额和短债同步扩张、文莱二期资本金与贷款条件缺乏约束，或剩余2.31亿库存股在用途不透明的情况下回流，治理风险将从折价因素升级为目标价下调因素。\textbf{行动上，出现上述组合时，即使利润仍高，也不提高估值倍数。}
\end{riskbox}

相反，若半年报和后续公告显示关联交易现金结算顺畅、担保与短债下降、员工计划费用透明、转债转股带来净债务下降，以及二期融资符合明确回报门槛，治理折价可以收窄。当前证据尚未达到这一升级条件。
